\section{Optional isolated Docker labs and Raspberry Pi deployment}
\subsection{Supported environments and actual validation}
For the requested Raspberry Pi OS machine with 8 GB RAM, confirm the architecture:
\begin{verbatim}
uname -m
getconf LONG_BIT
docker version
docker compose version
\end{verbatim}
The intended Pi configuration reports aarch64 and 64. An armv7l/32-bit OS is not a
validated target; use 64-bit Raspberry Pi OS for these profiles. Install Docker
Engine and the Compose plugin using the current Debian/Raspberry Pi OS instructions.
Build directly on the Pi unless you have a working multi-architecture build setup.

The pinned Python, Juice Shop, WebGoat, Mosquitto and HAProxy image indexes publish
AMD64 and ARM64 variants. The tests recorded here executed on an AMD64 laptop;
manifest availability does not establish Pi runtime performance. Budget one optional
lab at a time on the 8 GB Pi, monitor memory, and keep several GB free for the OS,
Docker builds and other services. The tracker is limited to 512 MB; Juice Shop to
768 MB; WebGoat to 1536 MB with a 768 MB maximum heap; MQTT to 128 MB; each gateway
to 64 MB. These are container limits, not measured worst-case RAM requirements.

The default tracker Compose file preserves its existing named volume and project
behavior. Do not change the project name on an existing installation without migrating
the volume. \texttt{ASCEND\_PORT} defaults to 8420. Read-only root, a bounded temporary
filesystem, UID 10001, dropped capabilities, process limits and bounded logs apply.
The pinned dependency constraints are in \texttt{docker/constraints.txt}.

\subsection{Start one optional lab}
The lab project is separate from the tracker. Starting it never mounts learner data
or the Docker socket in the vulnerable service. These commands run from the repo root:
\begin{verbatim}
docker compose -f docker/labs.compose.yml --profile juice-shop up -d --wait
docker compose -f docker/labs.compose.yml ps
# Stop the lab and remove its disposable state before another learner uses it:
docker compose -f docker/labs.compose.yml --profile juice-shop down
# Alternative, heavier Java application:
docker compose -f docker/labs.compose.yml --profile webgoat up -d --wait
# Alternative, deliberately anonymous teaching broker:
docker compose -f docker/labs.compose.yml --profile mqtt up -d --wait
\end{verbatim}
Use localhost:3000 for Juice Shop, localhost:8080/WebGoat for WebGoat,
localhost:9090/WebWolf for WebWolf, or localhost:1883 for MQTT. Profile port overrides
are \texttt{JUICE\_PORT}, \texttt{WEBGOAT\_PORT}, \texttt{WEBWOLF\_PORT}, and
\texttt{MQTT\_PORT}. Lab state is deliberately disposable; record notes first.
Different profiles may coexist, but running them together is not the Pi default.

Targets have separate internal Docker networks. Narrow HAProxy gateways join only
their assigned target network and a separate front network; they publish loopback
ports and fixed TCP forwards. Tests found that publishing ports directly from an
internal-only network on Docker 29.9.2 left them inaccessible; these gateways make
access explicit while keeping targets without public egress. No shared network
connects a vulnerable target to Ascend. Docker shares the host kernel; this is
containment for training, not a VM-grade boundary against kernel exploitation.

For a laptop browser connecting to labs on the Pi, use an SSH tunnel:
\begin{verbatim}
ssh -N -L 8420:127.0.0.1:8420 -L 3000:127.0.0.1:3000 YOUR_USER@lab-host.local
\end{verbatim}
Then open localhost on the laptop. Keep vulnerable services off public interfaces.
For phone access to Ascend itself, use the authenticated TLS deployment described in
the operations chapter. Do not reverse-proxy the vulnerable lab publicly alongside it.

WebGoat's pinned upstream image had a health check that invoked absent curl. The
replacement uses its included JDK and a mounted, fixed Java HTTP probe. Juice Shop
uses its included Node runtime for health checks. Some upstream challenges require
external blockchain/LLM services and cannot work on the isolated network. The bundled
verified tasks use local challenges only. WebGoat and MQTT are practice environments;
no automatic WebGoat or MQTT task-completion adapter is claimed in this release.

\subsection{Connect an external game to verified progress}
Initialize the private collector state on the Docker host before adding the overlay:
\begin{verbatim}
python3 scripts/collect_lab_proof.py init
docker compose -f docker/docker-compose.yml -f docker/lab-proofs.compose.yml up -d --build
python3 scripts/collect_lab_proof.py juice-shop --owner 1
\end{verbatim}
The collector writes \texttt{.ascend-labs/proofs.env} with mode 600 and signed
observations under its receipts directory. That private state is ignored by Git and
the build context. Do not print or commit the key. The overlay reads the private env
file and mounts only receipts read-only into Ascend. Mount creation fails if setup
was skipped. An absolute \texttt{ASCEND\_LAB\_STATE\_DIR} overrides its host path;
pass the same directory through the collector's \texttt{--state-dir} option.

For a different lab Compose project, use \texttt{--project PROJECT} on the collector.
For a different learner, use their numeric account ID (single mode is 1). The
administrator can list IDs without exposing credentials:
\begin{verbatim}
docker compose -f docker/docker-compose.yml exec ascend python -c \
 'from ascend.storage import Database; d=Database(); print([tuple(r) for r in d.connection.execute("SELECT id,username FROM users")]); d.close()'
\end{verbatim}
Each Juice Shop container must belong to one learner. The collector refuses to
reassign a running container; remove and recreate it for a new learner, and do not
share its gateway while a scored session runs. Concurrent collectors are serialized
on the host. This keeps assignment records coherent; it does not identify the person
at a shared browser. The game result is a bounded observation, not a proctored exam.

Refresh the collector after solving, then use \textbf{Verify lab completion} within
five minutes. The observation contains challenge identifiers, never a reusable
learner-submitted solved flag. Ascend checks its signature, owner, provider and age.
In public deployments collection is an administrator duty; the web server never
receives Docker access. The collector can be invoked manually after a solve; automated
polling is optional administration work and is not installed by this release.

\subsection{Bandit for a public starter pack}
Home hack lab v4 includes challenges 0--32 and preserves the old bandit-33 ID as an
endpoint/reflection task. The official level-34 page says it is not available; the
pack does not invent a further login. The series now starts at zero for new learners.
Existing completed IDs remain completed after an update and are not retroactively
certified by the new verifier.

Use the official initial login instructions, verify the official host key in your
SSH known-hosts file, solve a challenge, then run:
\begin{verbatim}
python3 scripts/collect_lab_proof.py bandit --owner 1 --level 0
ascend check home-hack-lab:bandit-0 --complete
\end{verbatim}
The collector authenticates to bandit1 on the fixed official SSH endpoint and checks a private short-lived SSH control connection. No remote command or forwarding
is requested, so restricted Bandit shells do not invalidate a successful login.
The collector closes the connection in cleanup. Password entry stays in the SSH terminal; passwords are not captured
or stored. Repeat with the challenge level just solved. One receipt records the
current check; verify before collecting the next level. Account access does not prove
that the puzzle was solved independently. Network availability and the game operator's
rules still apply. Automated tests stub SSH authentication validation and do not claim a
live Bandit login. Home lab host checks need actual SSH credentials and trust on the
machine running the verifier; the unprivileged Ascend container does not silently
inherit your host's SSH keys.

\subsection{Labs that need another environment}
GOAD requires Windows virtual machines; its Docker tooling does not make an AD domain
run as ARM Linux containers. Windows AD/AD CS and enterprise endpoint tasks need a
suitable VM host or hosted lab. Security Onion requires x86-64. FLARE-VM is a Windows
analysis environment; REMnux currently offers ARM64 appliances but is not bundled
here. Mobile tasks require the specified emulator/device; iOS work also has platform
requirements. CloudGoat uses real AWS resources and may incur charges; nothing here
creates a cloud tenant or deploys it. The OT/ICS track uses simulators and synthetic
traffic only. No live malware, radio exploitation or production control equipment
is launched by the Docker profiles.

\subsection{Reproduce the integration checks}
\begin{verbatim}
.venv/bin/python -m scripts.lab_smoke
.venv/bin/python scripts/container_smoke.py ascend:verified-labs
.venv/bin/python scripts/restore_smoke.py ascend:verified-labs
\end{verbatim}
The lab smoke test uses a disposable project and temporary learner database, validates
health and loopback routes, performs a real MQTT round trip, checks containment,
solves the Juice Shop scoreboard through a browser, and exercises the signed receipt
with correct and wrong accounts. It removes its own project afterward. Docker and
Playwright Chromium must be installed. Retained transcripts state the actual platform
and results. The archive/restore test uses disposable volumes; it never restores over
an existing learning profile.
